\section{Related Work}

\paragraph{VLM evaluation and robustness.} A large body of work measures VLM
accuracy on visual question answering and reasoning. Our focus is orthogonal:
not whether a model is correct on a fixed image, but whether it \emph{updates}
its decision under a controlled change. Citations are tracked in the
related-work matrix; no fabricated references are included.

\paragraph{Counterfactual and minimal-pair evaluation.} Counterfactual VQA and
minimal-pair probing construct paired inputs differing in one factor. CertVIC
differs by (i) using \emph{real} images edited under documented recipes rather
than synthetic or text-only counterfactuals, (ii) enforcing single-factor
validity through edit-quality gates and human review, and (iii) certifying the
resulting gap rather than reporting point estimates.

\paragraph{Image editing for evaluation.} Inpainting and object removal have been
used to create test stimuli. We treat edit realism and single-factor validity as
first-class measured quantities (quality gates plus inter-annotator agreement)
rather than assumed properties, and we keep crude and photorealistic edit engines
explicitly separated.

\paragraph{Anytime-valid inference.} Confidence sequences and betting-based tests
provide time-uniform coverage under optional stopping. We apply these tools to an
\emph{evaluation gap} so that a small, budget-constrained study yields a valid
certified lower bound; this statistical framing is the core methodological
contribution.

\paragraph{Dataset licensing and reproducibility.} Recipe-first release (pointers,
hashes, scripts) avoids redistributing non-rehostable pixels while remaining
reproducible.

Empirical comparisons remain [RESULT REQUIRED]. Full citations will be added
without fabrication.
